\chapter*{不定积分每日一题}

\refer{《积分的方法与技巧》金玉明\ 顾新身\ 毛瑞庭\ /\ 编著}

\exercise{求积分$\int{\cos{x}\qty(x^3+2x^2+3x+4)\d x}$}

\teacher{徐薪雨}

\postdate{2022年12月2日}

\result{\qty(x^3+2x^2-3x)\sin{x}+\qty(3x^2+4x-3)\cos{x}+C}

\analysis{
    本题考察多次分部积分法.


    分部积分的一般表达式是：
    \formula{\int{uv^{\prime}}\d x=uv-\int{v\d u}=uv-\int{vu^{\prime}\d x}}
    我们把分部积分的这一个过程成称为一升一降，一升的意思是$v^{\prime}$求原函数放到了$\d$后面变成了$v$，一降的意思是$u,\ v$交换位置后， 到了$\d$后面，再求导变成$\d u=u^{\prime}\d x$.于是，在我们的对应里，升阶意味着求原函数，降阶意味着求导.使用分部积分的重点，就在于 对被积函数进行划分，确定哪部分需要升阶（放到$\d$后面），哪部分需要降阶.而这样的对应只是便于利用直观辅助而已，不具备完全严谨的理由.在直觉中，降阶暗示着把一个函数由复杂形式变成简单形式；而升阶，往往意味着把函数变复杂.


下面来看今天的这个例题
\formula{\int{\cos{x}\qty(x^3+2x^2+3x+4)\d x}}
我们可以注意到，$\cos{x}$升多少次阶都不困难，而对$x^3+2x^2+3x+4$每降阶一次，就会次数减一，并且任意一个多项式经过有限次的降阶，最终都会得到$0$，而到$0$这一步时 就意味着答案全然浮现.


而这里是利用了多次分部积分，因此我们介绍一下分部积分的推广公式：


假定$u\qty(x)$和$v\qty(x)$在所考虑的区间上有直到$n+1$阶的连续导数：
\formula{u^{\prime}\quad u^{\prime\prime}\quad u^{\prime\prime\prime}\quad \cdots\quad u^{\qty(n)}\quad u^{\qty(n+1)}}
\formula{v^{\prime}\quad v^{\prime\prime}\quad v^{\prime\prime\prime}\quad \cdots\quad v^{\qty(n)}\quad v^{\qty(n+1)}}
那么通过不断的分部积分，不断的将$v$的导数升阶（放到$\d$后面），我们可以得到分部积分法的推广公式：
\formula{\int{uv^{\qty(n+1)}\d x}=uv^{\qty(n)}-u^{\prime}v^{\qty(n-1)}+u^{\prime\prime}v^{\qty(n-2)}-\cdots +\qty(-1)^nu^{\qty(n)}v+\qty(-1)^{n+1}\int{u^{\qty(n+1)}v\d x}}
那么回到上面的题目，我们可以用列表的方式去使用这个公式，令
\formula{u=x^3+2x^2+3x+4,\ \d v=\cos{x}\d x,\ v=\int{\cos{x}\d x}}
则
\formula{
    \addtolength{\jot}{-1em}
    \begin{aligned}
        & && v^{\qty(n+1)}=\cos{x}\\
        & u=x^3+2x^2+3x+4 && v^{\qty(n)}=\sin{x}\\
        & u^{\prime}=3x^2+4x+3 && v^{\qty(n-1)}=\cos{x}\\
        & u^{\prime\prime}=6x+4 && v^{\qty(n-2)}=-\sin{x}\\
        & u^{\prime\prime\prime}=6 && v^{\qty(n-3)}=-\cos{x}\\
        & u^{\qty(4)}=4 && v^{\qty(n-4)}=\sin{x}
    \end{aligned}
}
这样的话，我们就可以得到
\formula{
    \addtolength{\jot}{-1em}
    \begin{aligned}
        \int{\cos{x}\qty(x^3+2x^2+3x+4)\d x}=&\int{uv^{\qty(n+1)}\d x}\\
        =&uv^{\qty(n)}-u^{\prime}v^{\qty(n-1)}+u^{\prime\prime}v^{\qty(n-2)}-u^{\prime\prime\prime}v^{\qty(n-3)}\\
        =&\qty(x^3+2x^2+3x+4)\sin{x}-\qty(3x^3+4x+3)\qty(-\cos{x})\\
        &+\qty(6x+4)\qty(-\sin{x})-6\cos{x}+C\\
        =&\qty(x^3+2x^2-3x)\sin{x}+\qty(3x^2+4x-3)\cos{x}+C
    \end{aligned}
}
到这，这题就已经结束了.


稍微拓展一下，利用上述分部积分的方法，我们可以计算下面这些形式的积分：
\formula{\int{P\qty(x)\ee^{ax}\d x},\ \int{P\qty(x)\sin{bx}\d x},\ \int{P\qty(x)\cos{bx}\d x}}
其中$P\qty(x)$是$x$的多项式.操作方法与上题一样，感兴趣的同学可以尝试一下哦.
}